\providecommand{\td}{\operatorname{td}}
\providecommand{\Lie}{\operatorname{Lie}}
\providecommand{\Sym}{\operatorname{Sym}}
\section{The fixed-curve Picard construction}

All cohomology groups in this note have rational coefficients unless
specified otherwise. The divisor algebra of an abelian variety is the
subalgebra of its cohomology generated by rational divisor classes.\footnote{The construction obstruction concerns the fixed-curve Picard ansatz stated below. It does not assert that every possible algebraic correspondence is generated by divisors.}

\begin{theorem}\label{round26:picard:grr}
Let $A$ and $Z$ be complex abelian varieties, with $\dim A=a\ge2$, such that
\[
 \Hom_{\mathrm{HS}}(H^1(A),H^1(Z))=0.
\]
Let $C$ be a smooth projective connected curve with a chosen point,
$u:C\to A$ a morphism, and $v:Z\to\operatorname{Pic}^0(C)$ a homomorphism.
Pull the normalized Poincar\'e line bundle back to $C\times Z$ and call
it $\mathcal P$. For a vector bundle $F$ on $C$ of rank $r$, put
\[
 K=R(u\times1_Z)_*(\operatorname{pr}_C^*F\otimes\mathcal P)
       \in D^b(A\times Z).
\]
Write $z_A\in H^{2a}(A)$ for the class of a point and
\[
 D_v=\frac12(\operatorname{pr}_Z)_*c_1(\mathcal P)^2\in H^2(Z).
\]
Then $D_v$ is an algebraic divisor class and
\begin{equation}\label{round26:picard:character}
 \ch(K)=r\,u_*[C]\boxtimes1
        +\chi(C,F)z_A\boxtimes1
        +r\,z_A\boxtimes D_v.
\end{equation}
In particular the full character belongs to the divisor algebra of
$A\times Z$.

Let $\Phi_A$ be the relative Fourier--Mukai functor from $A\times Z$
to $\widehat A\times Z$ with normalized Poincar\'e kernel on
$A\times\widehat A$. There is a rational divisor class $D_u$ on
$\widehat A$ such that
\begin{equation}\label{round26:picard:fourier-character}
 \ch(\Phi_AK)
   =\chi(C,F)+r\,\operatorname{pr}_{\widehat A}^*D_u
                  +r\,\operatorname{pr}_Z^*D_v.
\end{equation}
Finite sums, shifts, cones and derived tensor products of these
transformed objects have characters in the divisor algebra.
\end{theorem}

\begin{proof}
Set $p=c_1(\mathcal P)$. Normalization and degree zero on the curve
fibres imply
\[
 p\in H^1(C)\otimes H^1(Z).
\]
Consequently $p^3=0$. If $\eta_C$ is the class of a point on $C$,
Riemann--Roch on the curve gives
\[
 \ch(F)\td(C)=r+\chi(C,F)\eta_C.
\]
The tangent bundles of $A$ and $Z$ are trivial. Applying
Grothendieck--Riemann--Roch to $q=u\times1_Z$, with the derived
pushforward in the statement, therefore gives
\begin{align*}
 \ch(K)
 &=q_*\bigl((r+\chi(C,F)\eta_C)\exp(p)\bigr)\\
 &=r q_*1+\chi(C,F)q_*\eta_C+r q_*p+\tfrac r2q_*p^2.
\end{align*}
All omitted products vanish by degree on $C$. The use of the derived
pushforward makes this identity valid without a weak-index or local
freeness hypothesis on the Picard transform. GRR in this form is the
smooth-target case of \cite[Chapter~18]{Ful98}.

The four displayed terms have K\"unneth bidegrees
\[
 (2a-2,0),\quad(2a,0),\quad(2a-1,1),\quad(2a,2).
\]
The third term is a rational Hodge class. Under Poincar\'e duality on
$A$, classes of type $(a,a)$ in
$H^{2a-1}(A)\otimes H^1(Z)$ are precisely morphisms of Hodge
structures $H^1(A)\to H^1(Z)$. Thus $q_*p=0$ by hypothesis.
The first two terms are respectively $u_*[C]\boxtimes1$ and
$z_A\boxtimes1$. Since the fourth has top degree on $A$, its
coefficient on $Z$ is obtained by integration over $A$ and equals
$(\operatorname{pr}_Z)_*p^2$. This proves \eqref{round26:picard:character}.
The class $D_v$ is algebraic directly from its formula, which is the
pushforward of an intersection of divisor classes. It lies in degree
two, so it is a divisor class.

To check the assertion about the divisor algebra, choose an ample
class $\theta_A$. Hard Lefschetz makes
\[
 \theta_A^{a-2}:H^2(A)\longrightarrow H^{2a-2}(A)
\]
an isomorphism of rational Hodge structures with its indicated shift.
The inverse image of $u_*[C]$ is of type $(1,1)$ and is therefore a
rational divisor class by Lefschetz $(1,1)$. Also $z_A$ is a rational
multiple of $\theta_A^a$. Every term of
\eqref{round26:picard:character} is consequently in the divisor algebra.

For the Fourier assertion, let $P_A$ be the normalized Poincar\'e
bundle on $A\times\widehat A$ and write $q_A=c_1(P_A)$. The
cohomological Fourier transform is
\[
 \mathcal F_A(\alpha)
  =\operatorname{pr}_{\widehat A,*}
      \bigl(\operatorname{pr}_A^*\alpha\exp(q_A)\bigr).
\]
This is the GRR action of the Fourier--Mukai functor because both
Todd classes are one, as in \cite[\S\S2,4]{Round26Mukai81}.
The class $q_A$ has bidegree $(1,1)$. On $H^j(A)$ only the term
$q_A^{2a-j}/(2a-j)!$ survives integration, so the output has degree
$2a-j$ on $\widehat A$. In particular
$\mathcal F_A(z_A)=1$ and
$D_u:=\mathcal F_A(u_*[C])\in H^2(\widehat A)$.
The latter is an algebraic class since the Fourier kernel is
algebraic, and hence a rational divisor class. Applying
$\mathcal F_A\otimes1$ to \eqref{round26:picard:character} gives
\eqref{round26:picard:fourier-character}. Additivity of the Chern character
in a distinguished triangle and its multiplicativity for derived
tensor products prove the final assertion.
\end{proof}

\begin{corollary}\label{round26:picard:mumford}
Let $X,Z$ be a full-group Mumford fourfold and the auxiliary sixfold
of the mixed-link construction. The fixed-curve Picard construction
of Theorem~\ref{round26:picard:grr}, with $A=X^2$, followed by the indicated
operations cannot have a Chern-character component projecting
nontrivially to the mixed adjoint link.
\end{corollary}

\begin{proof}
The representations $H^1(X)$ and $H^1(Z)$ have no common
irreducible constituent, so the hypothesis of
Theorem~\ref{round26:picard:grr} holds. Use a polarization to identify
$\widehat {X^2}$ with $X^2$ up to isogeny. Pullback along an isogeny
preserves the rational divisor algebra.
There are no cross divisor classes between $X^2$ and $Z$.
In the K\"unneth component
$H^1(X)\otimes H^1(X)\otimes H^2(Z)$, a product of two divisors
therefore has its first two factors in the one-dimensional
Poincar\'e pairing subspace. This is the trivial $G$-representation.
The rank-nine summand $T_X$ is a sum of adjoint representations
and has zero intersection with that subspace. Thus the mixed
projection of every divisor product, and hence of every character
in the stated construction, is zero.
\end{proof}

\begin{figure}[tb]
\centering
\begin{tikzpicture}[>=Latex,font=\small,node distance=10mm,
 every node/.style={align=center},box/.style={draw,rounded corners,fill=blue!6,text width=34mm,minimum height=12mm}]
\node[box] (c) {fixed curve $C$\\Poincar\'e twists};
\node[box,right=of c] (p) {Picard transform\\on $X^2\times Z$};
\node[box,right=of p] (d) {divisor-generated\\Chern character};
\draw[->] (c)--(p); \draw[->] (p)--(d);
\node[draw,fill=red!5,below=13mm of p,text width=65mm] (t) {projection to the exceptional mixed link is zero};
\draw[->] (d.south)|-(t.east);
\node[below=11mm of t,text width=120mm] {$\dim B\le5$: no single abelian carrier of $T_X$.\quad $\dim Z=6$: a Hodge carrier exists; its mixed algebraic cycle is the remaining construction.};
\end{tikzpicture}
\caption{Two restrictions on the mixed Mumford construction. The dimension bound concerns one abelian variety, as powers can contain additional representations.}
\label{round26:fig:hodge-carriers}
\end{figure}

\section{The smallest abelian carrier of the adjoint Hodge structure}

Let $F$ be a totally real cubic field and let $D/F$ be a quaternion
algebra split at one real place and ramified at the other two. Put
\[
 G=\Res_{F/\QQ}\SL_1(D),\qquad T_0=\Lie(G).
\]
Assume the split-corestriction condition so that the Mumford
representation $V$ exists over $\QQ$, with dimension eight and
\[
 G_{\overline\QQ}\cong\SL_2^3,\qquad
 V_{\overline\QQ}\cong U_1\otimes U_2\otimes U_3.
\]
Here $U_i$ is the standard representation at the $i$-th embedding.
These are the usual assumptions of Mumford's construction, as
reviewed in \cite{Galluzzi00}. We first give the rational
representation argument independently of Hodge theory.

\begin{lemma}\label{round26:carrier:small-representations}
Every irreducible rational representation of $G$ of dimension at
most ten is the trivial representation, $V$, or $T_0$.
Their dimensions are respectively one, eight and nine.
\end{lemma}

\begin{proof}
Representations of $G$ in characteristic zero are semisimple.
Over $\overline\QQ$, the irreducible representations are
\[
 W_{b_1b_2b_3}
  =\Sym^{b_1}U_1\otimes\Sym^{b_2}U_2\otimes\Sym^{b_3}U_3,
 \qquad b_i\ge0,
\]
of dimension $(b_1+1)(b_2+1)(b_3+1)$.
If $W$ is rationally irreducible, its distinct geometric
constituents form one Galois orbit, all with the same multiplicity.
Indeed the sum of the isotypic spaces in any Galois orbit descends
to a rational invariant subspace. More than one orbit would
contradict rational irreducibility. Equality of multiplicities
follows by applying Galois to the decomposition.

The Galois action on the three factors is the transitive
permutation action on the embeddings of $F$. The quaternion
forms do not change the highest-weight labels inside a factor,
because a root system of type $A_1$ has no diagram automorphism.
Every nonconstant triple consequently has an orbit of size at
least three. If its total orbit dimension is at most ten, its
individual dimension is at most three. The only nontrivial
possibilities are the orbits of $(1,0,0)$ and $(2,0,0)$, of total
dimensions six and nine before a possible common multiplicity.
For constant triples the only nontrivial possibility below
dimension eleven is $(1,1,1)$, of dimension eight. A constant
triple with entries at least two has dimension at least 27.

The six-dimensional possibility cannot descend with multiplicity
one to $\RR$. At a compact real factor, its complex standard
summand is the standard representation of $\mathrm{SU}(2)$,
with multiplicity one and with the other two factors acting
trivially. Real conjugation preserves this summand because its
complex representation type is invariant under the real form
and differs from the other two summands. A real structure would
therefore give an equivariant antilinear involution on
$\CC^2$. The standard $\mathrm{SU}(2)$-module has an
equivariant antilinear operator
\[
 J(z_1,z_2)=(-\overline z_2,\overline z_1),\qquad J^2=-1.
\]
By Schur's lemma every equivariant antilinear operator is $cJ$.
Its square is $-|c|^2$, so it cannot be an involution.
Thus a rational representation with that orbit has multiplicity
at least two and dimension at least twelve.

The remaining two geometric types are realized by $V$ and
$T_0$, whose split decompositions are respectively
$W_{111}$ and $W_{200}\oplus W_{020}\oplus W_{002}$.
They are rationally irreducible by the same Galois-orbit argument.
There is no second rational form of either type as a $G$-module.
For if $W$ has the same split type, then
\[
 \Hom_G(W,V)\otimes\overline\QQ
    =\Hom_{G_{\overline\QQ}}(W_{\overline\QQ},V_{\overline\QQ})
\]
is nonzero, and similarly with $T_0$. This base-change identity
holds because equivariance is a system of rational linear
equations on a finite-dimensional space of linear maps. A
nonzero rational equivariant map between the corresponding
irreducibles is an isomorphism. The only irreducible trivial
representation has dimension one. This completes the list.
\end{proof}

\begin{lemma}\label{round26:carrier:quotient-lift}
Let $H$ be a connected reductive group over $\QQ$, with a
surjection $H\to G^{\mathrm{ad}}$. There is a central isogeny
\[
 G\times H'\times S\longrightarrow H,
\]
where $H'$ is semisimple and $S$ is a torus. The restriction
of the surjection to $G$ is the simply connected covering of
$G^{\mathrm{ad}}$, after identifying its source with $G$.
For any rational representation of $H$, each rational
$G$-isotypic component is $H$-stable.
\end{lemma}

\begin{proof}
The connected centre of $H$ maps trivially to the adjoint
semisimple target. Thus $H^{\mathrm{der}}$ surjects onto
$G^{\mathrm{ad}}$. Take its simply connected cover. It is a
product of simply connected almost-simple groups over $\QQ$.
Their images are connected normal subgroups of
$G^{\mathrm{ad}}$. This target is $\QQ$-simple, because Galois
acts transitively on its three geometric simple factors.
Every nontrivial image is consequently the whole target.
Two different source factors cannot both have this image,
since their images commute and the target is noncommutative.
There is therefore exactly one such factor. Its kernel is
finite by almost-simplicity, so it is a simply connected cover
of $G^{\mathrm{ad}}$. Uniqueness of that cover identifies it
with $G$. The other factors form $H'$.

Multiplication of the derived group and the connected centre
is a central isogeny onto a connected reductive group. This
gives the displayed map with $S=Z(H)^\circ$. These standard
structure statements are recorded in \cite[Proposition~21.60,
Theorems~21.51 and~24.3]{Round26Milne17}.
On the source of this central isogeny, $H'$ and $S$ commute
with $G$. They preserve its isotypic components, and $G$
itself preserves them. Surjectivity proves $H$-stability.
\end{proof}

\begin{lemma}\label{round26:carrier:no-exterior-adjoint}
For $0\le j\le8$, the representation $\bigwedge^jV$ has no
constituent isomorphic to $T_0$ over $\QQ$.
\end{lemma}

\begin{proof}
It suffices to exclude each split adjoint
$A_i=\Sym^2U_i$. For odd $j$, the central element
$(-1,1,1)$ acts as $-1$ on $\bigwedge^jV$ and as $+1$ on
every $A_i$. For $j=0,8$ the exterior power is trivial.
In degree two,
\[
 \bigwedge^2V\cong\mathbf1\oplus
 (A_1\otimes A_2)\oplus(A_1\otimes A_3)
                         \oplus(A_2\otimes A_3).
\]
Degree six is the dual of degree two, since $\det V$ is
trivial. It remains to check degree four.

Its weights arise by choosing four distinct elements of the
eight sign triples. Let $w_{abc}$ denote their multiplicity.
The required counts are
\[
\begin{array}{c|rrrrrrrr}
(a,b,c)&200&400&220&202&420&402&222&422\\\hline
w_{abc}&4&1&2&2&0&0&1&0.
\end{array}
\]
For $200$, choose the one triple with negative first
coordinate. The omitted triple among the four with positive
first coordinate must have the same last two signs, giving
four choices. For $220$, those second signs are forced to be
positive for the negative-first triple and negative for the
omitted triple, leaving two choices of the third sign.
For $222$ both remaining signs are forced, giving one choice.
For $400$ the four positive-first triples are forced, and
their other coordinate sums are zero. This also gives all
three stated zero counts. The count for $202$ follows by
permuting coordinates.
Successive subtraction of raised weights, or multiplication
of the character by the three $\SL_2$ Weyl denominators,
gives the multiplicity of highest weight $(2,0,0)$ as
\[
 w_{200}-w_{400}-w_{220}-w_{202}
       +w_{420}+w_{402}+w_{222}-w_{422}=0.
\]
The other two adjoints have the same multiplicity by
permutation. This proves the assertion in every degree.
\end{proof}

Let $X$ now be a full-group Mumford fourfold for the preceding
data. Let $T_X\subset H^1(X)\otimes H^1(X)$ denote its
weight-two adjoint Hodge structure of dimension nine. Its
underlying $G$-module is $T_0$, and its Hodge group is
$G^{\mathrm{ad}}$.

\begin{theorem}\label{round26:carrier:minimal}
For every complex abelian variety $B$ of dimension at most
five and every integer $0\le r\le\dim B$,
\[
 \Hom_{\mathrm{HS}}
       \bigl(T_X,H^{2r}(B)(r-1)\bigr)=0.
\]
Consequently an abelian variety whose second cohomology
contains $T_X$ has dimension at least six.
\end{theorem}

\begin{proof}
Suppose there is a nonzero Hodge morphism of the indicated
form. The Hodge structure $T_X$ is simple: after splitting,
its three distinct adjoint constituents form a single
Galois orbit. Thus the map is injective.

Put $W=H^1(B)$ and $H=\Hg(W)$. It is connected reductive by
polarizability, see \cite[Proposition~3.6]{Del82}. Its action
on $\bigwedge^{2r}W$ preserves the image of $T_X$. Indeed a
Hodge subspace is preserved by the Hodge circle, and its
stabilizer is a rational algebraic subgroup, so it is
preserved by the rational Zariski closure of that circle.
The induced image of $H$ on this subspace is precisely its
Hodge group. Algebraic images commute here with taking the
Zariski closure of the defining circle. Tate twists do not
alter the action of the special Hodge circle. We obtain a
surjection $H\to\Hg(T_X)=G^{\mathrm{ad}}$.

Apply Lemma~\ref{round26:carrier:quotient-lift} and restrict $W$ to
its simply connected $G$-factor. Its dimension is at most
ten. Lemma~\ref{round26:carrier:small-representations} shows that
its only possible nontrivial rational irreducibles are $V$
and $T_0$. A copy of $T_0$ could occur only with multiplicity
one because its dimension is nine. Its rational isotypic
component would be $H$-stable by
Lemma~\ref{round26:carrier:quotient-lift}, hence a rational Hodge
substructure of $W$. This is impossible for a
nine-dimensional real space of weight one. The value of
the Hodge circle at $i$ preserves that space and squares to
$-1$, and a real vector space admitting a complex structure
has even dimension.

The $G$-action on $W$ is nontrivial. Otherwise the
$G$-factor would act trivially on the faithful defining
representation of $H$, contradicting its nontrivial image
in $H$. Thus a copy of $V$ occurs. It occurs once, and
\[
 W|_G\cong V\oplus\QQ^{\,d},\qquad d\le2.
\]
If $\dim W<8$, this already gives a contradiction.
For the remaining dimensions, the abelian cohomology
identity $H^*(B)=\bigwedge^*W$ shows that every $G$-type in
$H^{2r}(B)$ is a constituent of some $\bigwedge^jV$.
Lemma~\ref{round26:carrier:no-exterior-adjoint} excludes $T_0$
from all these representations. This contradicts the
assumed injection of $T_X$ and proves the theorem.
\end{proof}

\begin{remark}
The lower bound is attained at the level of Hodge structures
by the auxiliary sixfold of the mixed-link construction.
The theorem concerns the cohomology of one abelian variety.
Tensor powers of its cohomology can contain additional
representations, as $H^2(X^2)$ already contains $T_X$.
The Picard obstruction concerns a fixed curve and a fixed
bundle varied by Poincar\'e line bundles. These conclusions
leave the algebraic realization of the generic sixfold
mixed link as a separate problem.
\end{remark}
